\section{The two-face ten-nanometre specification}
\label{sec:accuracy}

Each Cartesian target is fixed in laboratory coordinates, with a declared
vertex, clear aperture and compatible non-optical continuation. The two
targets may be chosen independently as geometries; two arbitrary Cartesian
diopters do not automatically compose an aplanatic or stigmatic singlet.
The present objective is geometrical shaping. The optical quantities in
\cref{sec:optics} provide optional downstream diagnostics.

For mean graph surfaces $\bar h_j$ on apertures $D_{o,j}$, define
\begin{equation}
 e_{\rm mean}(t)=\max_{j=f,b}\sup_{D_{o,j}}|\bar h_j-h_{C,j}|,
 \qquad
 \epsilon=10^{-8}\ \mathrm m.
 \label{eq:mean-maximum-spec}
\end{equation}
No piston, tilt, best-fit sphere, or refocusing is removed. RMS error cannot
replace this maximum. The spatial search must include interior extrema between
nodes. A sampled maximum is a lower bound on the continuum supremum until an
interpolation or derivative bound encloses the omitted points.

With a single carrier and first-order normal displacement $\Xi_j$, the
instantaneous graph is
$h_j=\bar h_j+\Rea[(\Xi_j/n_{z,j})e^{-\mathrm i\omega t}]+O(|\Xi|^2)$.
At frozen slow time, maximization over the carrier phase and the aperture gives
\begin{equation}
 e_{\rm phase}=\max_{j=f,b}\sup_{D_{o,j}}
          \left(|\bar h_j-h_{C,j}|+\frac{|\Xi_j|}{|n_{z,j}|}\right).
 \label{eq:instantaneous-maximum-spec}
\end{equation}
This is exact for the retained single-harmonic graph model. Multiple carriers,
beat motion and nonlinear oscillations need their actual time-dependent sum
or a justified bound. Excluding optical index modulation from the objective
does not exclude physical motion of the surface.

Acceptance over a declared holding interval $[T_f,T_h]$ requires a justified
upper bound on the selected physical maximum throughout that interval. An error budget
must account for mean fit, numerical error, uncertainty, and
disturbances without treating unknown contributions as zero. Fast motion is
added only for the stricter instantaneous specification; its indirect effects
on the mean remain model-validity questions. Fixed-command
refinement supplies a numerical error estimate, not automatically a rigorous
upper bound. Model omission bounds require independent physics or measurements.
Formation starts at the solved unforced equilibrium with sources off; a flat
surface is a permissible initial condition only when it solves that problem.
Stored energy, fluid velocity and thermal state must also be initialized.

For both faces, the slow linearized block contains
\[
 \mathsf K_{\rm eff}=\mathsf K_{\rm cap+grav}
  -\begin{pmatrix}
       D_{h_f}q_f&D_{h_b}q_f\\D_{h_f}q_b&D_{h_b}q_b
    \end{pmatrix}.
\]
It is projected using the \emph{actual} volume constraints, coupled to the fluid
mass/mobility and linearized about the actual mean flow. The angular sectors
and memory requirements of \cref{sec:stability} apply to this block system.
Thus reaching a small static residual, maintaining a cycle-mean shape, and
maintaining the instantaneous physical tolerance are separately testable
scientific milestones.

\paragraph{Selected specification, 23 September 2026.}
The user-selected objective is the \emph{cycle-mean driven-liquid surface},
not the instantaneous interface and not a cured solid. Thus the active
geometric gate is $e_{\rm mean}\leq\epsilon$ on both apertures. Carrier
amplitude is not added to this error and need not itself be below ten
nanometres. Equation~\eqref{eq:instantaneous-maximum-spec} remains a separate,
stricter diagnostic, as do the earlier carrier-aware experiments.
The mean specification still needs a holding-interval error budget, numerical
convergence and mean-flow stability. Radiation traction, nonlinear averaging,
streaming and thermal feedback can change the cycle mean; removing the
instantaneous-motion penalty does not justify neglecting those mechanisms.
